\chapter{Bloco 4 -- Representação de valores discretos}
\label{cap:bloco4}

\begin{receita}[variável que só assume valores de uma lista]
Se $x$ só pode valer $v_1,\dots,v_k$ (por exemplo, $\{4,6,8,12,20,24\}$), use uma binária por valor:
\[
x=\sum_{i=1}^{k} v_i\,y_i,\qquad \sum_{i=1}^{k} y_i = 1,\qquad y_i\in\B.
\]
A segunda restrição escolhe \textbf{exatamente um} valor, e a primeira ``lê'' qual foi escolhido.

\textbf{Bônus:} qualquer função de $x$ também fica linear nos $y_i$:
\[
\phi(x)=\sum_{i=1}^{k}\phi(v_i)\,y_i,
\]
\textbf{mesmo que $\phi$ seja não linear}, porque $\phi(v_i)$ é só uma constante para cada $i$.
\end{receita}

\section{O problema: diâmetros de tubulação}

Três trechos de tubulação, com comprimentos $L_j$, precisam transportar uma vazão mínima $q_j$. Cada trecho usa
\textbf{um} diâmetro comercial. A vazão máxima de um tubo de diâmetro $d$ é $0{,}5\,d^2$ L/s (\textbf{não linear} em
$d$). O estoque de cada diâmetro é limitado. Queremos minimizar o custo total.

\begin{center}
\begin{tabular}{lcccccc}
\toprule
diâmetro $d$ (cm) & 4 & 6 & 8 & 12 & 20 & 24\\
\midrule
preço (R\$/m) & 8 & 12 & 18 & 30 & 70 & 100\\
vazão máx.\ $0{,}5d^2$ (L/s) & 8 & 18 & 32 & 72 & 200 & 288\\
estoque (m) & 500 & 500 & 500 & 150 & 250 & 300\\
\bottomrule
\end{tabular}
\qquad
\begin{tabular}{lccc}
\toprule
trecho $j$ & 1 & 2 & 3\\
\midrule
$L_j$ (m) & 300 & 200 & 250\\
$q_j$ (L/s) & 15 & 60 & 150\\
\bottomrule
\end{tabular}
\end{center}

\begin{discussao}
\begin{enumerate}
  \item \emph{Por que não uma variável inteira $d_j$ com $4\le d_j\le24$?} Porque 5, 7, 9, \ldots{} não existem
  comercialmente, e preço e vazão não são funções lineares de $d$.
  \item \emph{A vazão cresce com $d^2$; como o modelo continua linear?} Os coeficientes $0{,}5v_i^2$ são
  \textbf{constantes} (8, 18, 32, \ldots). O que varia são as binárias $y_{jd}$.
  \item \emph{Qual restrição acopla os trechos?} O \textbf{estoque}. Sem ele, cada trecho escolheria sozinho o
  diâmetro mais barato que dá vazão: $(6,12,20)$, custo 27\,100.
\end{enumerate}
\end{discussao}

\section{Construindo o modelo}

\textbf{Variáveis.} $y_{jd}\in\B$ = 1 se o trecho $j$ usa o diâmetro $d$ (3 trechos $\times$ 6 diâmetros = 18
binárias).

\textbf{Objetivo.} O custo de usar $d$ no trecho $j$ é $L_j\cdot\text{preço}_d$ (constante!):
\[
\min\sum_{j}\sum_{d} L_j\,\text{preço}_d\;y_{jd}.
\]

\textbf{Restrições.}
\[
\begin{aligned}
&\textstyle\sum_d y_{jd}=1 &&\forall j &&\text{(um diâmetro por trecho)}\\
&\textstyle\sum_d 0{,}5d^2\,y_{jd}\ge q_j &&\forall j &&\text{(vazão: linear nos }y\text{)}\\
&\textstyle\sum_j L_j\,y_{jd}\le \text{estoque}_d &&\forall d &&\text{(metros usados de cada diâmetro)}
\end{aligned}
\]

\begin{codigo}
def modelo_tubos():
    prob = pulp.LpProblem("tubos", pulp.LpMinimize)
    y = {(j, d): pulp.LpVariable(f"y_{j}_{d}", cat="Binary") for j in range(3) for d in D}
    # TODO 1: custo total
    prob += pulp.lpSum(L[j]*preco[d]*y[j, d] for j in range(3) for d in D), "custo"
    for j in range(3):
        # TODO 2: exatamente um diâmetro por trecho
        prob += pulp.lpSum(y[j, d] for d in D) == 1, f"um_diametro_{j}"
        # TODO 3: vazão mínima (vazao[d] = 0.5*d**2 é um NÚMERO)
        prob += pulp.lpSum(vazao[d]*y[j, d] for d in D) >= q[j], f"vazao_{j}"
    # TODO 4: estoque de cada diâmetro
    for d in D:
        prob += pulp.lpSum(L[j]*y[j, d] for j in range(3)) <= estoque[d], f"estoque_{d}"
    return prob, y
\end{codigo}

\section{Resultado e interpretação}

Para entender a resposta, liste as opções que dão vazão suficiente em cada trecho:
\begin{center}
\begin{tabular}{lll}
\toprule
trecho & diâmetros com vazão suficiente & custo de cada um\\
\midrule
1 ($q=15$, 300 m) & 6, 8, 12, 20, 24 & 3\,600; 5\,400; 9\,000; 21\,000; 30\,000\\
2 ($q=60$, 200 m) & 12, 20, 24 & 6\,000; 14\,000; 20\,000\\
3 ($q=150$, 250 m) & 20, 24 & 17\,500; 25\,000\\
\bottomrule
\end{tabular}
\end{center}
\begin{itemize}
  \item O trecho 2 queria o diâmetro 12, mas só há 150 m em estoque e ele precisa de 200 m. Então o trecho 2 usa 20
  ou 24.
  \item O diâmetro 20 tem 250 m: dá para o trecho 3 (250 m) \emph{ou} para o trecho 2 (200 m), não para os dois.
  \item Opções: trecho 2 com 20 e trecho 3 com 24 custa $14\,000+25\,000=39\,000$; trecho 2 com 24 e trecho 3 com 20
  custa $20\,000+17\,500=37\,500$. A segunda é melhor.
\end{itemize}
Ótimo: diâmetros $(6,\,24,\,20)$, custo $3\,600+20\,000+17\,500=\mathbf{41\,100}$.

\begin{atencao}[erros típicos]
\begin{itemize}
  \item Esquecer o estoque: dá 27\,100 (a dica da verificação avisa).
  \item Usar $\le 1$ em ``um diâmetro por trecho'': o modelo pode não escolher nenhum, mas aí a vazão falha;
  use $=1$, que é o que o problema diz.
  \item Escrever a vazão como \py{0.5*d**2} com \py{d} sendo uma variável do PuLP: não é linear. A potência é
  calculada \emph{sobre os dados}, antes, no dicionário \py{vazao}.
\end{itemize}
\end{atencao}

\begin{resumo}[Para pensar: variável semicontínua]
Uma variável $x\in\{0\}\cup[\ell,u]$ (``ou desligada, ou entre $\ell$ e $u$'') combina este bloco com o Bloco 1:
\[
\ell\,y\le x\le u\,y,\qquad y\in\B.
\]
Com $y=0$, $x=0$; com $y=1$, $\ell\le x\le u$.
\end{resumo}
